% !TeX program = xelatex
% ==========================================================================
% Пропоузал семестрового проєкту — курс «Штучний інтелект», УКУ.
% Це ЄДИНИЙ файл, який вам потрібно редагувати. Уся верстка й брендинг
% живуть у ucuproposal.cls — не змінюйте шрифти, поля, кеглі й кольори.
%
% Компіляція: XeLaTeX + Biber (Overleaf робить це автоматично).
% Локально:   latexmk (див. README.md).
% ==========================================================================
\documentclass[ukrainian]{ucuproposal}
% Опції класу (можна комбінувати через кому):
%   ukrainian | english  — мова документа й логотипу (за замовчуванням ukrainian)
%   slogan                — додає гасло УКУ в колонтитул титулки (за замовчуванням вимкнено)

% -------------------------------------
% Метадані титульної сторінки
% -------------------------------------
\projecttitle{Лічильник учасників дорожнього руху}


\teammember{Сиротюк Віктор}
\teammember{Британ Михайло}
\teammember{Мних Остап}
% \teammember{Прізвище Ім'я} % третій учасник (опційно)
% \teammember[https://linkedin.com/in/...]{Прізвище Ім'я} % опційне клікабельне посилання на профіль (LinkedIn тощо)

\mentor{Тарас Сливка} % якщо ментора немає — видаліть текст і залиште \mentor{}, рядок не друкується

\track{research} % research | engineering — не друкується на титулці, лише для внутрішньої категоризації

% \repo{https://github.com/team/project} % опційно; якщо репозиторію ще немає — залиште порожнім, рядок не друкується

\proposaldate{2026-09-20}

\begin{document}
\maketitlepage

\section{Проблема та мотивація}
% Що саме ви розв'язуєте і чому це важливо й цікаво в контексті ШІ?
% Орієнтовно 1 абзац.
Щоб розвивати міську інфраструктуру розумно, потрібні дані — хто ходить вулицею, коли і скільки. У багатьох європейських містах лічильники пішоходів і велосипедистів давно стали звичною річчю. В Україні таких систем майже немає, тому рішення про нові тротуари, переходи чи велосмуги часто приймаються без жодної фактичної бази. З погляду ШІ це класична, але
не тривіальна зв'язка детекції, багатооб'єктного трекінгу й підрахунку перетинів
лінії, яку треба втиснути в бюджет edge-пристрою та зробити стабільною в реальних
умовах: дощ, сутінки, перекриття, щільні групи людей. Окремий принциповий складник –
приватність: усе розпізнавання відбувається локально, а за межі пристрою виходять
лише знеособлені лічильники, без відео та зображень.

\section{Огляд джерел та наявних рішень}
% Що вже зроблено у світі за цією темою? Спирайтесь на references.bib,
% цитуйте через \cite{}. Орієнтовно 3-5 джерел.
Комерційні лічильники (Eco-Counter та подібні) розв'язують задачу «під ключ», але є
закритими й дорогими для масштабування на рівні міста~\cite{ecocounter}.
Дослідницька база натомість відкрита. Одноетапні детектори сімейства YOLO зробили
реальночасну детекцію придатною для вбудованих пристроїв~\cite{redmon2016,jocher2023}.
Для підрахунку потрібна не лише детекція, а й стійкі треки: SORT показав, що проста
зв'язка фільтра Калмана й угорського алгоритму дає високу швидкість~\cite{bewley2016},
DeepSORT додав ре-ідентифікацію за зовнішнім виглядом і зменшив кількість
перемикань ідентичності~\cite{wojke2017}, а ByteTrack продемонстрував, що врахування
детекцій з низькою впевненістю суттєво покращує трекінг у щільних сценах – саме
такі нас і цікавлять~\cite{zhang2022}. Для оцінювання спираємось на усталені метрики
багатооб'єктного трекінгу: CLEAR MOT~\cite{bernardin2008} і HOTA, яка збалансовано
враховує якість детекції та асоціації~\cite{luiten2021}. Перед стартом реалізації
плануємо окремо порівняти наявні open-source лічильники транспорту, щоб не
переписувати те, що вже працює.

\section{Дані}
% Які набори даних, звідки, ліцензія. Якщо збираєте самі — як саме.
% Якщо дані не потрібні — поясніть чому.
Ми не тренуємо детектор з нуля: базові ваги беремо з моделей, навчених на COCO~\cite{lin2014}, де
вже присутні класи \emph{person}, \emph{bicycle}, \emph{car}, \emph{bus},
\emph{truck}, \emph{motorcycle}. Для валідації й доналаштування збираємо власний
невеликий набір на пілотному майданчику – алея біля входу до Стрийського парку на
території УКУ. Записуємо кілька коротких сесій у різний час доби й погоду, анотуємо
вибірку кадрів (рамки + клас) і, окремо, ручні «еталонні» підрахунки за
15-хвилинними інтервалами – це наша ground truth для фінальної метрики. Орієнтовний
обсяг: 2-4 години відео, з них 1000-2000 анотованих кадрів. Публічні набори
багатооб'єктного трекінгу (MOT Challenge~\cite{milan2016}) використовуємо лише як
sanity-check пайплайну трекінгу, без донавчання на них – щоб не успадкувати їхні
некомерційні обмеження. Сирі відеозаписи
залишаються локально й не публікуються; відкритими робимо лише агреговані метадані
(тип, час, напрямок) під ліцензією CC~BY~4.0. Оскільки зйомка ведеться в публічному
просторі з телеоб'єктива, а на сервер не потрапляють ні кадри, ні обличчя, ні
номерні знаки, ризик для приватності мінімальний; перед збором узгодимо зйомку з
адміністрацією університету.
\section{Підхід і методи}
% Які алгоритми/архітектури плануєте використати, чи спираєтесь на
% існуючі реалізації та як плануєте їх покращити.
Пайплайн складається з чотирьох ланок. \textbf{(1)~Захоплення}: камера з
телеоб'єктивом, ручним фокусом і зумом, наведена на визначену ділянку з відстані
20-30~м. \textbf{(2)~Детекція}: попередньо натренована YOLO-модель у компактній
конфігурації~\cite{jocher2023}, квантизована під edge-інференс.
\textbf{(3)~Трекінг}: ByteTrack як основний варіант через його стійкість у щільних
сценах~\cite{zhang2022}, із SORT~\cite{bewley2016} як простішим baseline.
\textbf{(4)~Підрахунок}: віртуальні лінії та зони в кадрі; подія фіксується, коли
згладжена траєкторія треку перетинає лінію, і записується як
\texttt{(timestamp, class, direction, track\_id)}. Класи агрегуємо у дві групи –
вело-пішохідна зона (пішохід, велосипедист, користувач е-скутера) та проїжджа
частина (легковий автомобіль, вантажівка, громадський транспорт, мотоцикл); першу
версію свідомо звужуємо до пішоходів і велосипедистів, решту додаємо, якщо дозволить
час. Обчислення виконуються локально на міні-комп'ютері (рівня Raspberry~Pi~5 або
Jetson~Orin~Nano) поряд із камерою; назовні йдуть лише рядки метаданих. Наш внесок –
не нова архітектура, а адаптація: підбір геометрії ліній, фільтрація хибних і
подвійних спрацювань, доналаштування під конкретний ракурс і перевірка, що система
тримає прийнятний FPS на дешевому залізі.

\section{Оцінювання результатів}
% Метрики, baseline, спосіб порівняння, якісні очікування.
Головна метрика – точність підрахунку: відносна похибка (MAPE) між автоматичними й
ручними підрахунками за 15-хвилинними інтервалами, окремо для кожного класу й
напрямку. Цільове значення для пішоходів і велосипедистів – похибка до 10\,\% на
денних записах. Додатково міряємо якість класифікації (precision/recall, confusion
matrix між класами) і якість трекінгу за MOTA~\cite{bernardin2008} та
HOTA~\cite{luiten2021} на анотованій підмножині, а також продуктивність – FPS і
споживання пам'яті на цільовому пристрої. Baseline: (a)~детектор без трекінгу з
наївним підрахунком детекцій у зоні – показує, скільки дає саме асоціація;
(b)~зв'язка YOLO~+~SORT~\cite{bewley2016} проти основного варіанта з
ByteTrack~\cite{zhang2022}. Якісно очікуємо побачити деградацію в сутінках і при
щільних групах – ці режими виміряємо окремо й опишемо межі застосовності.

\section{План і ризики}
% Етапи до midterm і до фінального захисту, ключові ризики та план Б.
% Приклад врізки з ключовими цифрами (опційно, можна видалити):
% \begin{keyfacts}
% \textbf{Ключові цифри:} 3 датасети \textperiodcentered\ 2 baseline-моделі \textperiodcentered\ 4 тижні до midterm
\textbf{До midterm:} огляд аналогів і вибір обладнання; збір та анотація першого
відеонабору на пілотному майданчику; працездатний офлайн-прототип
«детекція~+~трекінг~+~підрахунок» на записаному відео; перші цифри похибки
підрахунку для пішоходів. \textbf{Після midterm:} перенесення інференсу на
міні-комп'ютер і оптимізація під реальний час; додавання класів велосипедистів і
транспорту; запис у сховище метаданих і простий дашборд; тривалий тестовий прогін у
полі; фінальний звіт із метриками та відкритим датасетом метаданих.

\textbf{Ризики та план Б.} Затримка чи брак заліза – розробляємо й оцінюємо все на
записаному відео, edge-деплой стає окремою, відокремлюваною фазою. Недостатній FPS
на пристрої – зменшуємо роздільну здатність і частоту кадрів, переходимо на меншу
квантизовану модель або обробляємо кожен n-й кадр. Погана точність у складних
умовах (ніч, дощ, щільний потік) – явно обмежуємо заявлену область застосовності
денним часом і доналаштовуємо модель на власних кадрах. Подвійний підрахунок через
розриви треків – додаємо гістерезис і мінімальну довжину траєкторії перед
зарахуванням події. Організаційний ризик із доступом до майданчика – маємо запасні
локації на території УКУ.

\begin{keyfacts}
\textbf{Ключові цифри:} 2-4 год власного відео \textperiodcentered\ 1000-2000 анотованих кадрів \textperiodcentered\ 2 baseline (без трекінгу, SORT) \textperiodcentered\ ціль: похибка підрахунку $\le$ 10\,\%
\end{keyfacts}

\section{Contribution statement}
% Хто за що відповідає. Відсотки не обов'язкові — якщо хочете їх вказати,
% додайте в квадратних дужках: \contribution[50\%]{Прізвище Ім'я}{...}
\begin{contributions}
  \contribution{Мних Остап}{детекція та класифікація: вибір і доналаштування моделі, квантизація, оцінювання якості}
  \contribution{Британ Михайло}{трекінг і логіка підрахунку, оптимізація та деплой на міні-комп'ютері, вимірювання продуктивності}
  \contribution{Сиротюк Віктор}{збір та анотація даних, ручні еталонні підрахунки, сховище метаданих і аналітика, документація}
\end{contributions}

% Примітка: декларація використання AI-інструментів (AI usage) до цього
% пропоузалу не входить — вона живе в README.md вашого репозиторію.

\ucuprintbibliography


\end{document}
